\chapter{Complete proofs for the sequential simulation targets}
\label{sec:seq-proof-completion}

This chapter completes the deductions used in the experimental chapters.
The foundational results and their proofs appear in
\Cref{sec:proof-foundations}; the sampling and interval proofs appear
in \Cref{sec:shared-statistical-methods}. A model assumption is an
input to these deductions. A recorded numerical observation is an output
of a specified finite computation, rather than a universal theorem.

\begin{proposition}[Finite and unbounded correlation domains]
\label{seqp:domain}
For a battery of size \(N\geq2\), the equicorrelation matrix
\(R_N=(1-c)I_N+c\mathbf1_N\mathbf1_N^{\mathsf T}\) is a covariance
matrix exactly when \(-1/(N-1)\leq c\leq1\). A fixed correlation is
admissible at arbitrarily large finite sizes exactly when \(0\leq c\leq1\).
The common-factor construction in \eqref{eq:seq-sim-factor} realizes every
finite size in this latter range.
\end{proposition}
\begin{proof}
Every vector is uniquely a multiple of \(\mathbf1_N\) plus an orthogonal
vector: subtract its coordinate average from each coordinate. On the first
subspace \(R_N\) multiplies by \(d_N=1+(N-1)c\), and on the second by
\(1-c\). The spectral covariance criterion in
Proposition~\ref{pf:spectral} therefore gives exactly the two inequalities
\(d_N\geq0\) and \(1-c\geq0\), including equality. Conversely, its
square-root construction supplies a Gaussian vector even at a zero
eigenvalue. A fixed \(c<0\) violates \(d_N\geq0\) once
\(N>1-1/c\); each \(0\leq c\leq1\) satisfies both inequalities for
every finite \(N\). Independent standard normals in
\eqref{eq:seq-sim-factor} give covariance \(c\) between distinct scores
and variance one, by Proposition~\ref{pf:moments}. Thus this construction
also proves existence in the range actually simulated. It makes no claim
to sample negative correlations, even though some finite negative
correlations are admissible.
\end{proof}

\begin{proposition}[The full conditional future law and point histories]
\label{seqp:conditioning}
For \(0\leq c<1\), the Gaussian residual construction gives a conditional
distribution at every specified initial vector \(z\). At
\(z=a\mathbf1_r\), it is given by the posterior factor formula
\eqref{eq:seq-sim-posterior-factor} and the entire future-vector
construction \eqref{eq:seq-sim-future-vector}. That future vector is jointly
Gaussian with mean
\(b_r\mathbf1_k\) and covariance
\begin{equation}
 K_k=(1-c)I_k+\frac{c(1-c)}{d_r}\mathbf1_k\mathbf1_k^{\mathsf T}.
 \label{eq:seq-complete-conditional-covariance}
\end{equation}
The mean function is the unique continuous conditional-mean version.
\end{proposition}
\begin{proof}
The initial covariance \(R_r\) is positive definite: its eigenvalues
are \(1-c>0\) and \(d_r>0\). Since
\((\mathbf1_r\mathbf1_r^{\mathsf T})^2
=r\mathbf1_r\mathbf1_r^{\mathsf T}\), multiplication gives
\[
 R_r^{-1}=\frac1{1-c}
   \left(I_r-\frac c{d_r}\mathbf1_r\mathbf1_r^{\mathsf T}\right),
 \qquad R_r^{-1}\mathbf1_r=\frac1{d_r}\mathbf1_r.
\]
The coefficient of \(\mathbf1_r\mathbf1_r^{\mathsf T}\) in the product
is \(-c/d_r+c/(1-c)-c^2r/[(1-c)d_r]=0\), because
\(d_r=1-c+rc\). This also verifies the formula when \(r=1\).

In this proof \(E_1,\ldots,E_k\) denote fresh future residuals,
distinct from every initial-block residual. Set
\[
 W=G-\frac{\sqrt c}{d_r}\sum_{i<r}X_i.
\]
The covariance identities of Proposition~\ref{pf:moments} give
\(\Cov(W,Z)=\sqrt c\mathbf1_r^{\mathsf T}
-\sqrt c\mathbf1_r^{\mathsf T}R_r/d_r=0\) and
\(\Var(W)=1-rc/d_r=(1-c)/d_r=\tau_r^2\).
The vector \((W,E_1,\ldots,E_k,Z)\) is Gaussian. The future residuals
are independent of \((G,Z)\), hence their covariances with \(W,Z\)
vanish; their covariance matrix is the identity. Gaussian block
independence (Proposition~\ref{pf:gaussian}) shows that
\((W,E_1,\ldots,E_k)\) is independent of \(Z\), and that its entries
are independent with variances \(\tau_r^2,1,\ldots,1\).
Consequently the conditional law of \(G\) at \(Z=z\) is represented by
\(\sqrt c\sum_i z_i/d_r+\tau_r H\), with a fresh standard normal
\(H\) independent of all future residuals. Substitution into the original
factor equation gives the whole future conditional vector, rather than
only its separate scalar moments. Its conditional covariance is
\((1-c)\delta_{ij}+c\tau_r^2\), proving the displayed matrix.

To check conditioning explicitly, write the future vector as
\(m(Z)+V\), with \(V\) independent of \(Z\). For any bounded Borel
\(h\) and Borel initial-history set \(A\), independence and Fubini give
\[
 \E[\mathbf1_{\{Z\in A\}}h(m(Z)+V)]
 =\int_A\!\int h(m(z)+v)\,P_V(dv)\,P_Z(dz).
\]
This is the defining conditional-kernel identity, as proved generally
in Proposition~\ref{pf:conditioning}. The affine mean is continuous;
the positive Gaussian density of \(Z\) implies that two continuous
versions agreeing almost everywhere agree everywhere: a nonzero
difference would persist on an open ball of positive probability.
Individual histories nevertheless have probability zero, so division by
\(P(Z=a\mathbf1_r)\) is not a conditional-probability calculation.
The displayed kernel specifies the version evaluated there. At \(c=1\),
all scores equal \(G\) almost surely, and the compatible equal-score
history gives future scores identically \(a\) directly on that support;
no singular covariance inverse is used.
\end{proof}

\begin{proposition}[Recursive predictions and the two aggregate laws]
\label{seqp:recursion}
For fixed equal initial values, all subsequent recursively filled
predictions are \(b_r=rca/d_r\). The fixed inclusive average is
\eqref{eq:seq-sim-centre} and has SD zero. The actual population and
conditional averages are Gaussian with the means and variances in
\eqref{eq:seq-sim-population}--\eqref{eq:seq-sim-conditional}.
\end{proposition}
\begin{proof}
For any prefix length \(s\), the inverse just proved and cross-covariance
\(c\mathbf1_s^{\mathsf T}\) give next-score mean
\(c\sum_{i<s}x_i/d_s\) by Proposition~\ref{pf:conditioning}. At the
initial history this is \(rca/d_r=b_r\). If the next \(j-1\) values
are all \(b_r\), the next value is
\[
 \frac c{d_r+(j-1)c}\bigl(ra+(j-1)b_r\bigr)
 =\frac c{d_r+(j-1)c}
   \frac{ra[d_r+(j-1)c]}{d_r}=b_r.
\]
Induction proves the assertion. Averaging gives
\((ra+kb_r)/N=r(d_r+kc)a/(Nd_r)=rd_Na/(Nd_r)\).
It is a deterministic number, whose centered square is zero, proving
SD zero. At \(c=1\), the compatible-support argument proves the same
conclusion with \(b_r=a\).

For the actual population sum, covariance bilinearity gives
\(N\) variance-one terms and \(N(N-1)\) ordered distinct-pair terms,
so its variance is \(N+N(N-1)c=Nd_N\). Its mean is zero. Division
by \(N^2\) proves the population formula. For the conditional sum, the
initial block is fixed. Summing every entry of the matrix in
\eqref{eq:seq-complete-conditional-covariance} gives
\[
 k(1-c)+k^2\frac{c(1-c)}{d_r}
 =\frac{k(1-c)(d_r+kc)}{d_r}
 =\frac{k(1-c)d_N}{d_r}.
\]
Dividing by \(N^2\) and summing the means gives the conditional formula.
Both are affine transformations of the relevant Gaussian vectors, so
they are Gaussian by Proposition~\ref{pf:gaussian}. Gaussianity is
essential here: matching two moments alone would not establish a normal
law. Finally, Proposition~\ref{pf:moments} gives
\(\E[100+15C]=100+15\E[C]\) and
\(\Var(100+15C)=225\Var(C)\), proving every reporting-scale conversion.
\end{proof}

\begin{proposition}[Population spread, conditional spread and random histories]
\label{seqp:variance-decomposition}
For \(0\leq c<1\), an unconstrained random initial block gives
\begin{align*}
 \E[C_{r,k}\mid Z]&=\frac{d_N}{Nd_r}\sum_{i<r}X_i,\\
 \Var(\E[C_{r,k}\mid Z])&=\frac{r d_N^2}{N^2d_r},\\
 \Var(C_{r,k})&=\frac{r d_N^2}{N^2d_r}
                 +\frac{k(1-c)d_N}{N^2d_r}=\frac{d_N}{N}.
\end{align*}
Thus conditional and population SD describe different experiments;
conditioning does not change the stipulated marginal population SD of
each original test.
\end{proposition}
\begin{proof}
Each future conditional mean at a general \(z\) is
\(c\sum_i z_i/d_r\) by Proposition~\ref{seqp:conditioning}. The first
identity follows by adding the fixed initial sum to those \(k\) means.
There are \(r\) diagonal variance terms and \(r(r-1)\) ordered
covariance terms in the initial sum, giving variance \(rd_r\); squaring
the coefficient proves the second identity. The residual conditional
variance is constant in \(z\), by the proved covariance matrix.
Total variance (Proposition~\ref{pf:expectation}) gives the sum. Expanding
\(r d_N+k(1-c)=Nd_r\) simplifies it as displayed. At the simulated
\(c=3/5\), all factors in the between-history variance are positive,
so conditioning strictly reduces variance. The conditional law's centre
changes because the observed block is fixed high, while the unconditioned
centre remains zero. For \(r=2,k=10\), multiplication by 225 yields the
exact IQ-scale decomposition
\(285/2=1805/16+475/16\). Taking square roots gives the reported SDs;
the decomposition adds variances, not SDs.
\end{proof}

\section{Exact certificates for the numerical target columns}
\label{seqp:exact-grid}
At \(c=3/5\), \(a=10/3\), and \(N=r+k\), the preceding proposition
gives the following rational IQ-scale targets:
\[
 m=100+\frac{50r(3N+2)}{N(3r+2)},\qquad
 V_P=\frac{45(3N+2)}N,\qquad
 V_C=\frac{90k(3N+2)}{N^2(3r+2)}.
\]
Indeed, \(d_j=(3j+2)/5\); inserting this into the proved centre and
variance formulas and multiplying variances by 225 gives the three
expressions. Direct substitution gives all finite exact entries below.
For \(k=\infty\), the proved limits in \Cref{sec:seq-sim-limits}
give \(m=100+150r/(3r+2)\), \(V_P=135\), and
\(V_C=270/(3r+2)\).

\begin{center}\small
\begin{tabular}{rr r r r}
\toprule
\(r\)&\(k\)&\(m\)&\(V_P\)&\(V_C\)\\\midrule
1&1&\(140\)&\(180\)&\(36\)\\
1&2&\(410/3\)&\(165\)&\(44\)\\
1&10&\(1450/11\)&\(1575/11\)&\(6300/121\)\\
2&1&\(875/6\)&\(165\)&\(55/4\)\\
2&2&\(575/4\)&\(315/2\)&\(315/16\)\\
2&10&\(1675/12\)&\(285/2\)&\(475/16\)\\
3&1&\(1625/11\)&\(315/2\)&\(315/44\)\\
3&2&\(1610/11\)&\(153\)&\(612/55\)\\
3&10&\(20450/143\)&\(1845/13\)&\(36900/1859\)\\\midrule
1&\(\infty\)&\(130\)&\(135\)&\(54\)\\
2&\(\infty\)&\(275/2\)&\(135\)&\(135/4\)\\
3&\(\infty\)&\(1550/11\)&\(135\)&\(270/11\)\\\bottomrule
\end{tabular}
\end{center}

Here is an exact rounding proof, independent of floating-point evaluation.
For \(m=u/v\), \(v>0\), a displayed \(t/T\) differs from \(m\)
by less than half a unit precisely when
\(D=v-2|Tu-tv|>0\). For an SD \(\sqrt{p/q}\), \(p,q>0\), and
positive displayed integer \(t\), the two certificates
\[
 L=4T^2p-(2t-1)^2q>0,\qquad
 U=(2t+1)^2q-4T^2p>0
\]
are equivalent to
\((t-1/2)/T<\sqrt{p/q}<(t+1/2)/T\): multiply by \(4T^2q\)
and square the positive endpoints. Thus positivity proves the nearest
rounding, with no tie convention needed. The following exact integers
certify all four-decimal target mean and SD entries in the tables, using
\(T=10^4\), the fractions above, and their displayed integers \(t\).
The population mean is exactly 100 in every row.

\begin{center}\scriptsize
\setlength{\tabcolsep}{4pt}
\begin{tabular}{rr r rr rr}
\toprule
\(r\)&\(k\)&\(D\)&\(L_P\)&\(U_P\)&\(L_C\)&\(U_C\)\\\midrule
1&1&1&621071&452241&239999&240001\\
1&2&1&848591&179025&528431&2225\\
1&10&7&1108621&9421371&24721751&45126225\\
2&1&2&848591&179025&584316&602276\\
2&2&4&1011982&996002&558704&5120784\\
2&10&4&1649950&260018&5138544&1835664\\
3&1&5&1011982&996002&9206676&211436\\
3&2&3&661775&327769&982625&13694895\\
3&10&123&12046427&343197&95794525&566797691\\
1&\(\infty\)&1&359&929161&113039&474841\\
2&\(\infty\)&2&359&929161&465116&1393924\\
3&\(\infty\)&9&359&929161&3790525&569259\\\bottomrule
\end{tabular}
\end{center}
Every listed integer is positive; the established equivalences therefore
prove the claimed roundings. For the six-decimal worked example in
\Cref{sec:seq-sim-sds}, the corresponding SD pairs \((L,U)\),
with \(T=10^6\), are \((169283518,21713858)\) for
\(\sqrt{285/2}\approx11.937336\) to the stated precision and
\((125135856,572288016)\) for \(\sqrt{475/16}\approx5.448624\).
For \(1675/12\) displayed as \(139.583333\), the mean certificate is
\(D=4\). These are strict rounding statements; the decimals do not equal
the irrational SDs exactly.

The theoretical MC SE columns use \(\sqrt{V/M}\) with
\(M=4\,000\,000\), so the same certificate applies with denominator
\(Mq\). In finite-grid row order above, their four-decimal display
integers \(t\) are
\[
 \begin{aligned}
  \text{Population:}\quad &(67,64,60,64,63,60,63,62,60),\\
  \text{Conditional:}\quad &(30,33,36,19,22,27,13,17,22).
 \end{aligned}
\]
Substitution of the listed exact variances into the two integer expressions gives
\(\min L=96\,000\,000\) and \(\min U=356\,000\,000\) across
these 18 entries. Both minima are positive, which proves every displayed
MC SE rounding by the same inequalities. Observed estimates, the
standardized discrepancies calculated from them, and confidence-interval
endpoints calculated from observed variances remain finite numerical
records; their reference-law interpretation is proved separately.

\section{What a proof of a recorded calculation means}
\label{seqp:record-scope}
The target columns and their count/limit comparisons are mathematical
deductions above. The observed columns, extrema, replicate comparisons,
CSV coordinates, graph sizes, seeds, software versions and historical
proof-check counts are finite reproduction records. Their evidential
claim is that the named source, inputs and reduction formulas produced
the recorded outputs; it is checked by reproducing that calculation and
comparing the record. A theorem about ideal Gaussian samples does not
deduce a particular seeded output. Conversely, agreement of finitely many
outputs cannot prove an identity for all model inputs. For example, a
function altered only at an untested input has exactly the same tested
outputs. This counterexample explains the logical scope of finite tests.

There are \(3\cdot3\cdot2=18\) finite laws and \(3\cdot2=6\)
limiting laws, hence 24 output cases; four one-million-observation
replicates give four million per case and 96 million scalar-mean output
observations. Reusing normal arrays across cases does not alter that
output count, but does prevent interpreting the cases as independent
replications. Pooling formulas are proved in
Proposition~\ref{sm:pooling}; ideal sample laws, interval inversion and
the scope of multiple comparisons are proved in
Proposition~\ref{sm:normal-sample}, Corollary~\ref{sm:sd-interval} and
Proposition~\ref{sm:multiple}. No event-count denominator is used in this sequential
experiment: every generated scalar mean contributes to its row. The
probability-zero initial history is handled by the conditional kernel,
not by accepting or rejecting population observations. Automated software
obligation counts concern the recorded source and its contracts; they
are neither proofs of this written derivation nor empirical validation of
the Gaussian model.
